\documentclass[11pt]{amsart}
\usepackage[margin=1.05in]{geometry}
\usepackage{amsmath,amssymb,amsthm,mathtools,xcolor,hyperref,enumitem}\sloppy
\hypersetup{colorlinks=true,linkcolor=blue!50!black,citecolor=blue!50!black,urlcolor=blue!50!black}

\newcommand{\srccommit}{PENDING}
\newcommand{\gr}[1]{\textsf{#1}}
\newcommand{\Hom}{\operatorname{Hom}}
\newcommand{\End}{\operatorname{End}}
\newcommand{\Ind}{\operatorname{Ind}}
\newcommand{\Res}{\operatorname{Res}}
\newcommand{\rad}{\operatorname{rad}}
\newcommand{\topp}{\operatorname{top}}
\newcommand{\Des}{\operatorname{Des}}
\newcommand{\SYT}{\operatorname{SYT}}
\newcommand{\kk}{\Bbbk}
\newcommand{\sq}{\boxtimes}
\newcommand{\ok}{\mathrm{ok}}
\newcommand{\Ib}{\mathcal I}
\newcommand{\Jb}{\mathcal J}

\theoremstyle{plain}
\newtheorem{theorem}{Theorem}
\renewcommand{\thetheorem}{\Alph{theorem}}
\newtheorem{corollary}[theorem]{Corollary}
\newtheorem{lemma}{Lemma}[section]
\newtheorem{proposition}[lemma]{Proposition}
\newtheorem{fact}[lemma]{Fact}
\newtheorem{lemmaII}{Lemma}
\renewcommand{\thelemmaII}{II.\arabic{lemmaII}}
\newtheorem{corII}[lemmaII]{Corollary}
\theoremstyle{definition}
\newtheorem{remark}[lemma]{Remark}
\newtheorem{remarkII}[lemmaII]{Remark}

\title[Two indecomposability theorems for $0$-Hecke modules]{Two indecomposability theorems for $0$-Hecke modules:\\ induced simples (Q72(i)) and Kim--Lee--Oh Conjecture 7.5 (Q73)}
\author{Rick}
\date{2026-11-11}

\begin{document}
\maketitle
\thispagestyle{empty}

\begin{center}\small
From: Rick. \quad To: Clio. \quad Date: 2026-11-11.\\
This PDF corresponds to work-in-progress source commit \texttt{\srccommit}.\\
Proof records (local): \texttt{proofs/2026-11-10-q72-ind-simples-indecomposable.md} (PROVE 289),\\
\texttt{proofs/2026-11-11-klo-conj75-indecomposable.md} (PROVE 290).
\end{center}

\medskip
\noindent\fbox{\parbox{\dimexpr\linewidth-2\fboxsep-2\fboxrule}{\small
\textbf{Summary.} Two results on induction products of $0$-Hecke modules, proved by one skeleton (Frobenius reciprocity for $\End$ and for the top, the Mackey double-coset splitting, and an induction on size).
\begin{itemize}[leftmargin=1.2em,itemsep=1pt]
\item \textbf{Theorem~\ref{thm:A} (Q72(i)).} $\Ind_{H_a(0)\otimes H_b(0)}^{H_{a+b}(0)}(C_I\sq C_J)$ is indecomposable unless $I,J\in\{(m),(1^m)\}$, in which case it is $P_{I\cdot J}\oplus P_{I\odot J}$.
\item \textbf{Theorem~\ref{thm:B} (Q73 = Kim--Lee--Oh, arXiv:2310.20571, Conj.~7.5).} If a skew shape is disconnected and no two consecutive connected components are ribbons, then $X_{\lambda/\mu}$ is indecomposable. Strong form: Corollary~\ref{cor:C}.
\end{itemize}
\textbf{Grades:} both \gr{proved (internal grade)}; each has passed one internal cold read (Dream 288 for Q72(i), Dream 289 for Q73), and the nits found there are fixed in this write-up. Neither has had an external check. Q73 relies on cited results of KLO, listed in \S\ref{sec:cited}.

\textbf{What I want from you:} a referee-style check of both proofs, especially \textbf{Lemma~\ref{II:key}} (Key Lemma and the return map $\zeta$; Lemma 3 in the PROVE 290 record) and \textbf{Lemma~\ref{II:support}} (support lemma; Lemma 5 in the record) of Part~II.}}

\tableofcontents

\section{Conventions and the shared word model}\label{sec:conv}

Throughout, $\kk$ is a field. $H_N=H_N(0)$ is the $\kk$-algebra generated by $\pi_1,\dots,\pi_{N-1}$ with $\pi_i^2=\pi_i$ and the braid relations. For $w\in S_N$ put $\pi_w:=\pi_{i_1}\cdots\pi_{i_\ell}$ for any reduced word $s_{i_1}\cdots s_{i_\ell}$ of $w$. The basic multiplication rule is
\begin{equation}\label{eq:mult}
\pi_i\pi_w=\begin{cases}\pi_{s_iw}, & \ell(s_iw)>\ell(w),\\ \pi_w, & \ell(s_iw)<\ell(w).\end{cases}
\end{equation}
For integers $p\le q$ we write $W_{[p,q]}:=\langle s_p,\dots,s_q\rangle$. This is a \emph{generator interval}: $W_{[p,q]}$ is the symmetric group on the values $\{p,\dots,q+1\}$ and fixes all other values. For $Z\subseteq[N-1]$, $W_Z$ is the parabolic subgroup generated by $\{s_z:z\in Z\}$.

\subsection{Simples, compositions, projectives}
For $D\subseteq[N-1]$ let $C_D$ be the one-dimensional module on which $\pi_i$ acts by $[i\in D]$. If $w\in S_N$, then $\pi_w$ acts on $C_Z$ by $[w\in W_Z]$, because the letters of any reduced word of $w$ form its support. A composition $I\vDash m$ is identified with its descent (partial-sum) set $D_I\subseteq[m-1]$, and we write $C_I:=C_{D_I}$. For $I\vDash a$ and $J\vDash b$, the concatenation $I\cdot J$ and near-concatenation $I\odot J$ have
\[
D_{I\cdot J}=D_I\cup\{a\}\cup(a+D_J),\qquad D_{I\odot J}=D_I\cup(a+D_J).
\]
$P_D$ (or $P_I$) denotes the projective cover of $C_D$.

\begin{fact}[Norton 1979]\label{fact:N}
The $C_D$, $D\subseteq[m-1]$, are all the simple $H_m$-modules, and $\dim P_D=\#\{w\in S_m:\Des(w)=D\}$. Hence $C_D$ is projective iff $D\in\{\emptyset,[m-1]\}$, i.e.\ iff $I\in\{(m),(1^m)\}$. ``Projective'' for a composition always means this.
\end{fact}

\subsection{The word model of an induced module}
Fix $a,b\ge1$, $N=a+b$, and let $B:=H_a\otimes H_b\subseteq H_N$ be the subalgebra generated by the $\pi_i$ with $i\ne a$; here $\pi_{a+q}$ is the $q$-th generator of $H_b$. For a $B$-module $V$ put $\Ind V:=H_N\otimes_B V$, and for an $H_a$-module $X$ and an $H_b$-module $Y$ put $X\sq Y:=\Ind(X\otimes Y)$.

Let $W^{(a,b)}$ be the set of \emph{shuffles}: the $w\in S_N$ that are increasing on $[1,a]$ and on $[a+1,N]$. These are the minimal-length left coset representatives of $S_a\times S_b$. $H_N$ is free as a right $B$-module on $\{T_w:=\pi_w : w\in W^{(a,b)}\}$, so $\Ind V=\bigoplus_{w\in W^{(a,b)}}T_w\otimes V$.

We encode a shuffle $w$ by a word $x(w)=x_1\cdots x_N\in\{A,B\}^N$ indexed by \emph{values}: $x_v=A$ iff $w^{-1}(v)\le a$. Every word with exactly $a$ letters $A$ comes from exactly one shuffle.

\begin{lemma}[action rules]\label{lem:rules}
Let $w\in W^{(a,b)}$, $v\in V$, and $1\le i\le N-1$. According to the letters $(x_i,x_{i+1})$ of $x(w)$:
\begin{enumerate}[label=\textup{(\roman*)},itemsep=1pt]
\item $AB$: $\pi_i(T_w\otimes v)=T_{s_iw}\otimes v$, and the word of $s_iw$ has $BA$ in positions $i,i+1$;
\item $BA$: $\pi_i(T_w\otimes v)=T_w\otimes v$;
\item $AA$: $\pi_i(T_w\otimes v)=T_w\otimes\pi_pv$, where $p$ is the rank of $i$ among the $A$-values;
\item $BB$: $\pi_i(T_w\otimes v)=T_w\otimes\pi_{a+q}v$, where $q$ is the rank of $i$ among the $B$-values.
\end{enumerate}
\end{lemma}

\begin{proof}
(i) Here $w^{-1}(i)\le a<w^{-1}(i+1)$, so $\ell(s_iw)=\ell(w)+1$. The permutation $s_iw$ exchanges two values lying in different blocks, so it is still increasing on both blocks, i.e.\ a shuffle. By \eqref{eq:mult}, $\pi_iT_w=T_{s_iw}$.
(ii) Here $w^{-1}(i)>w^{-1}(i+1)$, so $s_i$ is a left descent of $w$, and $\pi_iT_w=T_w$ by \eqref{eq:mult}.
(iii) The values $i,i+1$ are consecutive $A$-values. Since $w$ is increasing on $[1,a]$, they occupy consecutive positions $p,p+1$, where $p$ is the rank of $i$ among the $A$-values. Then $s_iw=ws_p$ with $\ell(ws_p)=\ell(w)+1$ (this is Deodhar's lemma in this case), so $\pi_iT_w=T_{ws_p}=T_w\pi_p$. Since $p\le a-1$, $\pi_p\in B$.
(iv) is the same as (iii), with positions $a+q,a+q+1$.
\end{proof}

\begin{lemma}[Mackey splitting]\label{lem:mackey}
For a shuffle $w$ let $k(w):=\#\{v\le a: x_v=B\}$, so $0\le k(w)\le\min(a,b)$. Put $M_k:=\bigoplus_{k(w)=k}T_w\otimes V$. Then $\Res_B\Ind V=\bigoplus_{k}M_k$ as $B$-modules, and $M_0=T_{\mathrm{id}}\otimes V\cong V$.
\end{lemma}

\begin{proof}
For $i\neq a$, each rule of Lemma~\ref{lem:rules} either keeps the word or exchanges $x_i$ and $x_{i+1}$, and $i,i+1$ lie on the same side of the cut between $a$ and $a+1$. So $k(w)$ is preserved. The only word with $k=0$ is $A^aB^b$, which is the word of $w=\mathrm{id}$.
\end{proof}

\begin{remark}\label{rem:mackey-any-q}
Lemma~\ref{lem:mackey} is the standard double-coset (Mackey) decomposition of $\Res_B\Ind_B$: the $M_k$ are indexed by the $(S_a\times S_b)$-double cosets. It is a direct sum for the Hecke algebra at every $q$, not just at $q=0$. What $q=0$ adds is the \emph{shape} of the pieces, recorded in Lemma~\ref{lem:cut}: every basis vector is sent by each $\pi_i$ to another basis vector or to a scalar multiple of an element of $V$ in the same slot, and this makes $M_k$ a tensor product of two smaller induced modules of the same type.
\end{remark}

\begin{lemma}[cutting a Mackey block]\label{lem:cut}
Let $k\ge1$. Cut each word as $x=(\ell\mid r)$ with $\ell=x_1\cdots x_a$ (containing $a-k$ letters $A$ and $k$ letters $B$) and $r=x_{a+1}\cdots x_N$ (containing $k$ letters $A$ and $b-k$ letters $B$). Every pair $(\ell,r)$ of such words occurs. On $T_{(\ell\mid r)}\otimes v$:
\begin{enumerate}[label=\textup{(\alph*)},itemsep=1pt]
\item for $i<a$, $\pi_i$ changes only $\ell$. In case $AA$ it acts on $v$ by $\pi_p$, where $p\le a-k-1$ is the rank of $i$ among the $A$'s \emph{of $\ell$}. In case $BB$ it acts by $\pi_{a+q}$, where $q\le k-1$ is the rank of $i$ among the $B$'s of $\ell$;
\item for $i=a+t$ with $1\le t\le b-1$, $\pi_i$ changes only $r$. In case $AA$ it acts by $\pi_{(a-k)+p'}$, where $p'\in[1,k-1]$ is the rank of $t$ among the $A$'s of $r$. In case $BB$ it acts by $\pi_{a+k+q'}$, where $q'\in[1,b-k-1]$ is the rank of $t$ among the $B$'s of $r$.
\end{enumerate}
\end{lemma}

\begin{proof}
Any word with $a$ letters $A$ is a shuffle word, so every pair occurs. Every value in $[1,a]$ is below every value in $[a+1,N]$. Hence the rank of a left letter among all $A$-values (resp.\ $B$-values) equals its rank inside $\ell$. A right $A$-letter has overall $A$-rank $(a-k)+{}$(its rank in $r$), and a right $B$-letter has overall $B$-rank $k+{}$(its rank in $r$). Now apply Lemma~\ref{lem:rules}. The bounds hold because both letters $i$ and $i+1$ are counted.
\end{proof}

\part*{Part I. Induced simples: Q72(i)}

\section{Statement}

Fix $I\vDash a$ and $J\vDash b$ with $a,b\ge1$, $N=a+b$. Let $V=C_I\sq C_J=\kk v$ as a $B$-module, so $\pi_i$ ($i\ne a$) acts on $v$ by $\chi(i):=[i\in D]$, where
\[
D:=D_I\cup(a+D_J)\subseteq[N-1]\setminus\{a\}.
\]
Put $M=M_{I,J}:=H_N\otimes_BV$ and $g:=1\otimes v$. By \S\ref{sec:conv}, $M$ has basis $e_w:=T_w\otimes v$ ($w\in W^{(a,b)}$). In Lemma~\ref{lem:rules}(iii),(iv) the actions become scalars: $\pi_ie_w=\chi(p)e_w$ in case $AA$ and $\pi_ie_w=\chi(a+q)e_w$ in case $BB$.

\begin{theorem}[Q72(i)]\label{thm:A}
$M_{I,J}$ is indecomposable unless $I$ and $J$ are both projective, i.e.\ $I,J\in\{(m),(1^m)\}$. In the both-projective case $M_{I,J}\cong P_{I\cdot J}\oplus P_{I\odot J}$, with two nonzero, non-isomorphic summands.
\end{theorem}

\section{\texorpdfstring{$\End M$}{End M} and the locality criterion}

\begin{lemma}[Frobenius]\label{lem:I-frob}
$\varphi\mapsto\varphi(g)$ is a linear isomorphism $\End_{H_N}(M)\to E:=\{u\in M:\pi_iu=\chi(i)u\ \text{for all } i\ne a\}$. Write $\varphi_u$ for the endomorphism with $\varphi_u(g)=u$.
\end{lemma}

\begin{proof}
$\End_{H_N}(H_N\otimes_BV,M)\cong\Hom_B(V,\Res M)$, and since $V=\kk v$ with $\pi_i v=\chi(i)v$, a $B$-map $V\to M$ is the same thing as a choice of image $u\in E$.
\end{proof}

\begin{lemma}[top]\label{lem:I-top}
$\Hom_{H_N}(M,C_Z)\cong\Hom_B(V,\Res C_Z)$. This is one-dimensional if $Z\setminus\{a\}=D$ and zero otherwise. Hence $\topp M=C_{Z_1}\oplus C_{Z_2}$ with $Z_1=D=D_{I\odot J}$ and $Z_2=D\cup\{a\}=D_{I\cdot J}$, and these are distinct since $a\notin D$. Let $f_j:M\to C_{Z_j}$ be the unique module maps with $f_j(g)=1$. Then
\[
f_j(e_w)=[w\in W_{Z_j}].
\]
\end{lemma}

\begin{proof}
The first isomorphism is Frobenius reciprocity. $\Res C_Z$ is the character $i\mapsto[i\in Z]$ ($i\ne a$), which equals $\chi$ iff $Z\setminus\{a\}=D$. Simples are one-dimensional with $\End=\kk$, so the multiplicity of $C_Z$ in $\topp M$ is $\dim\Hom(M,C_Z)$. Finally, $f_j(e_w)=f_j(\pi_wg)=\pi_w\cdot1=[w\in W_{Z_j}]$ in $C_{Z_j}$.
\end{proof}

\begin{proposition}[(Crit)]\label{prop:crit}
$M$ is indecomposable if and only if $f_1(u)=f_2(u)$ for every $u\in E$.
\end{proposition}

\begin{proof}
Each $\varphi\in\End M$ induces an endomorphism of $\topp M=C_{Z_1}\oplus C_{Z_2}$. Since the top is multiplicity-free, this endomorphism is $\operatorname{diag}(c_1,c_2)$, where $f_j\circ\varphi=c_jf_j$, so $c_j=f_j(\varphi(g))$. This defines an algebra map $\rho:\End M\to\kk\times\kk$.

The kernel $\ker\rho$ is $\{\varphi:\varphi(M)\subseteq\rad M\}$. Every such $\varphi$ is nilpotent, because $\varphi(\rad^jM)=\rad^j\varphi(M)\subseteq\rad^{j+1}M$. Hence $\ker\rho$ is a nil ideal, and $\ker\rho\subseteq J(\End M)$.

The image $\operatorname{im}\rho$ is a subalgebra of $\kk\times\kk$ containing $(1,1)$, so it is either $\kk(1,1)$ or $\kk\times\kk$.
\begin{itemize}[itemsep=1pt]
\item If $\operatorname{im}\rho=\kk\times\kk$, the idempotent $(1,0)$ lifts modulo the nil ideal $\ker\rho$ to a nontrivial idempotent of $\End M$, so $M$ is decomposable.
\item If $\operatorname{im}\rho=\kk(1,1)$, then $\End M/\ker\rho\cong\kk$ with $\ker\rho$ nil, so $\End M$ is local and $M$ is indecomposable.
\end{itemize}
Finally, $\operatorname{im}\rho=\kk(1,1)$ iff $c_1=c_2$ for all $\varphi$. By Lemma~\ref{lem:I-frob} this means $f_1(u)=f_2(u)$ for all $u\in E$.
\end{proof}

\section{Computing \texorpdfstring{$f_1$ and $f_2$}{f1 and f2}}

\begin{lemma}\label{lem:f1f2}
\textup{(a)} $f_1(e_w)=[w=\mathrm{id}]$, so $f_1(u)=u_{\mathrm{id}}$ (the coefficient of $e_{\mathrm{id}}=g$).

\textup{(b)} Let $[\alpha,\beta]\ni a$ be the component of $Z_2$ containing $a$ (as a generator interval). Then $\alpha=a-r'$ and $\beta=a+s'$, where $r':=$ the length of the run $a-1,a-2,\dots$ in $D_I$ and $s':=$ the length of the run $1,2,\dots$ in $D_J$. For a shuffle $w$,
\begin{equation}\label{eq:f2}
f_2(e_w)=[x_1=\dots=x_{\alpha-1}=A]\cdot[x_{\beta+2}=\dots=x_N=B].
\end{equation}
\end{lemma}

\begin{proof}
(a) Since $a\notin Z_1$, we have $W_{Z_1}\subseteq S_a\times S_b$, and the only shuffle in $S_a\times S_b$ is $\mathrm{id}$.

(b) The description of $[\alpha,\beta]$ is immediate from $Z_2=D_I\cup\{a\}\cup(a+D_J)$. Write $W_{Z_2}$ as the commuting product of the $W_K$ over the components $K$ of $Z_2$, and write $w\in W_{Z_2}$ as $w=\prod_Kw_K$. The right descent set of $w$ is the union of those of the $w_K$, and $w_K\ne e$ has a right descent in $K$. A shuffle has right descents only at $a$. So every factor $w_K$ with $a\notin K$ is trivial, and a shuffle lies in $W_{Z_2}$ iff it lies in $W_{[\alpha,\beta]}$, i.e.\ iff it fixes every value $<\alpha$ and every value $>\beta+1$.

Now let $w$ be a shuffle. If $w$ fixes $1,\dots,\alpha-1$, then these values sit in positions $\le\alpha-1\le a$, so they are $A$-values. Conversely, if $1,\dots,\alpha-1$ are $A$-values, then, since $w$ is increasing on $[1,a]$, they are the $\alpha-1$ smallest $A$-values and sit in positions $1,\dots,\alpha-1$. Symmetrically, using $\beta+2\ge a+2$, $w$ fixes every value $>\beta+1$ iff all those values are $B$-values.
\end{proof}

\section{Mackey factorisation of \texorpdfstring{$E$}{E}}

By Lemma~\ref{lem:mackey}, $\Res_BM=\bigoplus_kM_k$ with $M_0=\kk g$. Since the $\pi_i$ ($i\ne a$) preserve each $M_k$, we get $E=\bigoplus_k(E\cap M_k)$.

\begin{lemma}\label{lem:LR}
For $k\ge1$, the map $e_{(\ell\mid r)}\mapsto e_\ell\otimes e_r$ is a $B$-module isomorphism $M_k\cong L_k\otimes R_k$, where
\begin{align*}
L_k&:=\Ind_{H_{a-k}\otimes H_k}^{H_a}\bigl(C_{I^{(1)}}\sq C_{J^{(1)}}\bigr), & R_k&:=\Ind_{H_k\otimes H_{b-k}}^{H_b}\bigl(C_{I^{(2)}}\sq C_{J^{(2)}}\bigr),\\
D_{I^{(1)}}&=D_I\cap[1,a-k-1], & D_{I^{(2)}}&=\{p-(a-k):p\in D_I,\ p>a-k\},\\
D_{J^{(1)}}&=D_J\cap[1,k-1], & D_{J^{(2)}}&=\{q-k:q\in D_J,\ q>k\}.
\end{align*}
Degenerate cases: $L_a=C_{J^{(1)}}$ (word $B^a$) and $R_b=C_{I^{(2)}}$ (word $A^b$).
\end{lemma}

\begin{proof}
Compare Lemma~\ref{lem:cut} with the word models of $L_k$ and $R_k$. On the left, $AA$ at $i<a$ gives the eigenvalue $\chi(p)=[p\in D_I]$ with $p\le a-k-1$, i.e.\ $[p\in D_{I^{(1)}}]$. $BB$ gives $\chi(a+q)=[q\in D_J]$ with $q\le k-1$, i.e.\ $[q\in D_{J^{(1)}}]$. On the right, $AA$ gives $[(a-k)+p'\in D_I]=[p'\in D_{I^{(2)}}]$, and $BB$ gives $[k+q'\in D_J]=[q'\in D_{J^{(2)}}]$. The $AB$/$BA$ rules agree verbatim. $H_a$ acts only on $\ell$ and $H_b$ only on $r$, and the basis of $M_k$ is the product of the two sets of words.
\end{proof}

\noindent\textbf{Eigenspaces multiply.} Let $U_k:=\{x\in L_k:\pi_ix=[i\in D_I]x\ \forall i<a\}\cong\Hom(C_I,L_k)$ and $U'_k:=\{y\in R_k:\pi_ty=[t\in D_J]y\ \forall t<b\}\cong\Hom(C_J,R_k)$. Then
\[
E\cap M_k=U_k\otimes U'_k.
\]
Indeed, $\bigcap_i\ker((A_i-c_i)\otimes1)=(\bigcap_i\ker(A_i-c_i))\otimes R_k$ for commuting operators $A_i$ on $L_k$, likewise on the right, and $(U\otimes R_k)\cap(L_k\otimes U')=U\otimes U'$.

\noindent\textbf{$f_2$ on a block.} By \eqref{eq:f2}, and since $\alpha-1\le a-1$ and $\beta+2\ge a+2$, we get $f_2|_{M_k}=\lambda_k\otimes\lambda'_k$, where
\[
\lambda_k(e_\ell)=[\ell_1=\dots=\ell_{\alpha-1}=A],\qquad \lambda'_k(e_r)=[r_{\beta+2-a}=\dots=r_b=B].
\]
$f_2$ is an $H_N$-map to $C_{Z_2}$, and $\Res_BC_{Z_2}=C_I\sq C_J$. So $f_2|_{M_k}$ is a $B$-map $L_k\otimes R_k\to C_I\otimes C_J$. If $\lambda'_k\ne0$, choose $y$ with $\lambda'_k(y)=1$. Then $\lambda_k(\pi_ix)=(\lambda_k\otimes\lambda'_k)(\pi_ix\otimes y)=[i\in D_I]\lambda_k(x)$, so $\lambda_k:L_k\to C_I$ is a module map. Symmetrically, if $\lambda_k\ne0$ then $\lambda'_k:R_k\to C_J$ is a module map.

\begin{proposition}\label{prop:witness}
If $M$ is decomposable, then there is $k\ge1$ with $\lambda_k(U_k)\ne0$ and $\lambda'_k(U'_k)\ne0$. For such $k$:
\begin{enumerate}[label=\textup{(\alph*)},itemsep=1pt]
\item $C_I$ is a direct summand of $L_k$, and $C_J$ is a direct summand of $R_k$;
\item $k\le r'+1$.
\end{enumerate}
\end{proposition}

\begin{proof}
On $M_0=\kk g$ we have $f_1=f_2$, and for $k\ge1$, $f_1$ vanishes on $M_k$ by Lemma~\ref{lem:f1f2}(a). So for $u=\sum_ku_k\in E$ with $u_k\in E\cap M_k$, $f_2(u)-f_1(u)=\sum_{k\ge1}f_2(u_k)$. If $M$ is decomposable, Proposition~\ref{prop:crit} gives $u\in E$ with $f_1(u)\ne f_2(u)$, hence some $k\ge1$ with $(\lambda_k\otimes\lambda'_k)(U_k\otimes U'_k)\ne0$.

(a) Both $\lambda_k$ and $\lambda'_k$ are nonzero, so both are module maps. An $x\in U_k$ with $\lambda_k(x)\ne0$ spans a copy of $C_I$, and the composite $C_I\to L_k\xrightarrow{\lambda_k}C_I$ is nonzero, hence an isomorphism. So $C_I$ splits off $L_k$. The same argument works for $C_J$ and $R_k$.

(b) $\lambda_k(e_\ell)\ne0$ requires all $k$ letters $B$ of $\ell$ to lie in positions $[\alpha,a]$, which has $a-\alpha+1=r'+1$ slots.
\end{proof}

\section{Induction on \texorpdfstring{$N$}{N} and the converse}

We induct on $N$. The base case is $N=2$, i.e.\ $a=b=1$, where both simples are projective. Assume Theorem~\ref{thm:A} holds for all sizes $<N$, suppose $M_{I,J}$ is decomposable, and take $k$ as in Proposition~\ref{prop:witness}.

\begin{lemma}\label{lem:I-ind}
If $1\le k\le a-1$ and $C_I$ is a summand of $L_k$, then $I$ is projective. Likewise, if $1\le k\le b-1$ and $C_J$ is a summand of $R_k$, then $J$ is projective.
\end{lemma}

\begin{proof}
$L_k=M_{I^{(1)},J^{(1)}}$ is a module of size $a<N$, so the induction hypothesis applies. If $L_k$ is indecomposable, then $L_k\cong C_I$, so $\dim L_k=\binom ak=1$. That is impossible for $1\le k\le a-1$. Otherwise $L_k$ is decomposable, so by induction $I^{(1)},J^{(1)}$ are projective and $L_k\cong P\oplus P'$ is projective. Then $C_I$, as a summand of a projective, is projective. The right-hand statement has the same proof.
\end{proof}

Recall $k\le\min(a,b)$. We split into cases.
\begin{itemize}[itemsep=2pt]
\item $k<a$ and $k<b$. Lemma~\ref{lem:I-ind} applies on both sides, so $I$ and $J$ are projective.
\item $k=a<b$. $J$ is projective by Lemma~\ref{lem:I-ind}. On the left, $C_I\mid L_a=C_{J^{(1)}}$ gives $D_I=D_J\cap[1,a-1]$. Since $D_J\in\{\emptyset,[b-1]\}$, this is $\emptyset$ or $[a-1]$, so $I$ is projective.
\item $k=b<a$. Symmetrically, $I$ is projective and $D_J=D_{I^{(2)}}\in\{\emptyset,[b-1]\}$.
\item $k=a=b$. Here $L_a=C_{J^{(1)}}$ with $D_{J^{(1)}}=D_J$, so $D_I=D_J$. Also, Proposition~\ref{prop:witness}(b) gives $a\le r'+1$, so $r'=a-1$ and $[1,a-1]\subseteq D_I$. Hence $I=J=(1^a)$, which is projective.
\end{itemize}
So decomposability forces both $I$ and $J$ to be projective.

\noindent\textbf{Converse.} If $I$ and $J$ are projective, then $C_I=P_I$ and $C_J=P_J$. $P_I\otimes P_J$ is a projective $B$-module, and $H_N$ is free over $B$, so $M=\Ind(P_I\otimes P_J)$ is projective. The multiplicity of $P_Z$ in $M$ is $\dim\Hom(M,C_Z)$, which by Lemma~\ref{lem:I-top} is $1$ exactly for $Z\in\{Z_1,Z_2\}$. So $M\cong P_{I\cdot J}\oplus P_{I\odot J}$. \qed

\section{Remarks on Part I}

\begin{remark}[where $q=0$ is used]\label{rem:I-q0}
The proof applies Frobenius reciprocity twice: once to $\End M$ and once to $\topp M$. The Mackey decomposition $\Res_BM=\bigoplus_kM_k$ is the double-coset decomposition. It is a direct sum at \emph{every} $q$ (Remark~\ref{rem:mackey-any-q}), so its directness is not where $q=0$ enters. The proof uses $q=0$ in three places:
\begin{itemize}[itemsep=1pt]
\item the specific shape of the pieces: $M_k\cong L_k\otimes R_k$ with explicitly matched eigenvalues (Lemma~\ref{lem:LR}), which rests on the $0$-Hecke rules of Lemma~\ref{lem:rules};
\item the fact that the simples are one-dimensional characters $C_D$;
\item Norton's projectivity criterion (Fact~\ref{fact:N}).
\end{itemize}
Mackey then reduces the gluing question to ``does $C_I$ split off a smaller induced module $L_k$?'', and induction answers it.
\end{remark}

\begin{remark}
The bound $k\le r'+1$ is used only in the case $k=a=b$. It is what excludes, for example, $I=J=(2,2)$, where $\dim\End M=4$ but $M$ is indecomposable.
\end{remark}

\begin{remark}[negative control]
In the both-projective case a witness $k$ does exist. For example, take $k=1$ with $a,b\ge2$: then $L_1$ and $R_1$ are projective and contain $C_I$ and $C_J$. In the non-projective case, the hypothesis that fails is ``$C_I$ is a summand of $L_k$''. The Wake 288 guess (``$I$ has both an ascent and a descent'') is the $k<a$ shadow of this condition, via Fact~\ref{fact:N}.
\end{remark}

\begin{remark}
Hypotheses used: $q=0$ (as in Remark~\ref{rem:I-q0}), the two-block parabolic, and one-dimensional factors. Q72(ii), which asks for a description of $\End M$ in general, is untouched. The proof does describe $\End M$ modulo its radical: $\End M=\kk\oplus\ker\rho$.
\end{remark}

\part*{Part II. Kim--Lee--Oh Conjecture 7.5: Q73}

\section{Statement}

KLO is Kim--Lee--Oh, \emph{Regular Schur labeled skew shape posets and their $0$-Hecke modules}, arXiv:2310.20571. For a basic skew shape $\lambda/\mu$ with $n$ cells, $X_{\lambda/\mu}$ is the $H_n$-module with basis $\SYT(\lambda/\mu)$ (KLO App.~A). For an SYT $T$ and $1\le i\le n-1$:
\[
\pi_iT=\begin{cases}T,& i\text{ strictly left of }i+1,\\ 0,& i\text{ and }i+1\text{ in the same column},\\ s_iT,& i\text{ strictly right of }i+1.\end{cases}
\]
KLO show that $X_{\lambda/\mu}\cong M_P$ whenever $\operatorname{sh}(\tau_P)=\lambda/\mu$. Write $\lambda/\mu=C_1\star\cdots\star C_m$ for its connected components, with $C_1$ the north-east one. The components occupy pairwise disjoint row ranges and column ranges, and $C_s$ lies strictly north-east of $C_t$ for $s<t$. A component is a \emph{ribbon} if it contains no $2\times2$ square.

\begin{theorem}[KLO Conjecture 7.5]\label{thm:B}
If $m\ge2$ and no two consecutive components $C_s,C_{s+1}$ are both ribbons, then $X_{\lambda/\mu}$ is indecomposable.
\end{theorem}

\begin{corollary}[strong form]\label{cor:C}
For an arbitrary skew shape $\lambda/\mu$, $X_{\lambda/\mu}\cong\bigoplus_SX_S$ with every $X_S$ indecomposable. Here $S$ runs over the shapes obtained by replacing each maximal run of $r$ consecutive ribbon components by one connected ribbon, using one of the $2^{r-1}$ choices of $\cdot/\odot$ concatenation (KLO Lemma 5.1, Rem.~A.1(2)). So the number of indecomposable summands is $\prod2^{r-1}$.
\end{corollary}

\section{Word model for shapes, rotation}

The notation of \S\ref{sec:conv} is in force, with $X$ an $H_a$-module, $Y$ an $H_b$-module and $V=X\otimes Y$. In Lemma~\ref{lem:rules}(iv), $\pi_{a+q}$ acts on $Y$ as $\pi_q$.

\begin{lemmaII}[tableau $=$ word model; KLO Prop.~A.2(1), JKLO arXiv:2105.14169 Lemma~4]\label{II:W}
Let $S,S'$ be basic skew shapes with $|S|=a$ and $|S'|=b$. The map
\[
T_w\otimes T\otimes T'\ \longmapsto\ \text{the filling of }S\star S'\text{ by }w(T)\text{ on }S\text{ and }w(a+T')\text{ on }S'
\]
is an isomorphism $X_S\sq X_{S'}\to X_{S\star S'}$.
\end{lemmaII}

\begin{proof}
\emph{Bijection on bases.} There are no order relations between $S$-cells and $S'$-cells in $S\star S'$. So an SYT of $S\star S'$ is exactly a choice of the value set on $S$ (equivalently a shuffle $w$) together with two standardised SYT $T,T'$. The filling is standard because $w$ is increasing on both blocks.

\emph{Action.} Every $S$-cell is strictly to the right of every $S'$-cell, and in a different row. If $i\in S$ and $i+1\in S'$, the tableau rule gives $s_i$, which matches word rule $AB$. If $i\in S'$ and $i+1\in S$, it gives the identity, matching $BA$. If $i,i+1\in S$, they are the images of $p,p+1$ with $p$ the rank of $i$, and their relative position is that of $p,p+1$ in $T$. So the tableau rule is the $X_S$ rule for $\pi_p$, matching $AA$. Rule $BB$ is the same on the $S'$ side.
\end{proof}

\begin{lemmaII}[rotation]\label{II:rot}
Let $T^\circ$ be $T$ rotated by $180^\circ$, with each value $v$ replaced by $n+1-v$. Then $T\mapsto T^\circ$ gives $X_{(\lambda/\mu)^\circ}\cong\varphi^*X_{\lambda/\mu}$, where $\varphi$ is the automorphism $\pi_i\mapsto\pi_{n-i}$. Moreover $(C_1\star C_2)^\circ=C_2^\circ\star C_1^\circ$, and rotation sends ribbons to ribbons. In particular, rotation preserves indecomposability and the hypothesis of Theorem~\ref{thm:B}.
\end{lemmaII}

\begin{proof}
Rotation reverses left and right and preserves columns. So ``$i$ is strictly left of / in the same column as / strictly right of $i+1$ in $T^\circ$'' is equivalent to ``$n-i$ is strictly left of / in the same column as / strictly right of $n-i+1$ in $T$''.
\end{proof}

\section{Tops and the indecomposability criterion}

\noindent\textbf{Standing hypothesis (H).} $X$ and $Y$ are indecomposable, with multiplicity-free tops $\topp X=\bigoplus_{i\in\Ib}C_{D_i}$ and $\topp Y=\bigoplus_{j\in\Jb}C_{E_j}$. Let $\varepsilon_i:X\to C_{D_i}$ and $\delta_j:Y\to C_{E_j}$ be the top functionals, and fix $x_i\in X$ and $y_j\in Y$ with $\varepsilon_i(x_i)=1$ and $\delta_j(y_j)=1$.

\begin{lemmaII}[top]\label{II:top}
Assume only that $\topp X$ and $\topp Y$ are multiplicity-free. Then $\topp(X\sq Y)=\bigoplus_{i,j,e}C_{Z_{ije}}$ is multiplicity-free, where $e\in\{0,1\}$ and
\[
Z_{ij0}=D_i\cup(a+E_j),\qquad Z_{ij1}=Z_{ij0}\cup\{a\}.
\]
The corresponding functional is $f_{ije}(T_w\otimes v)=[w\in W_{Z_{ije}}]\cdot(\varepsilon_i\otimes\delta_j)(v)$.
\end{lemmaII}

\begin{proof}
Frobenius reciprocity and the tensor identity for finite-dimensional modules give
\[
\Hom(X\sq Y,C_Z)\cong\Hom_B(X\otimes Y,\Res C_Z)=\Hom(X,C_{Z\cap[a-1]})\otimes\Hom(Y,C_{(Z-a)\cap[b-1]}).
\]
By multiplicity-freeness this is one-dimensional exactly when $Z=Z_{ije}$ for some $i,j,e$, and zero otherwise. The extension of $\beta\in\Hom_B(V,\Res C_Z)$ is $T_w\otimes v\mapsto\chi_Z(\pi_w)\beta(v)=[w\in W_Z]\beta(v)$. The sets $Z_{ije}$ are pairwise distinct, since $Z$ determines $D_i$, $E_j$ and whether $a\in Z$.
\end{proof}

\begin{lemmaII}[criterion]\label{II:crit}
Assume (H) and put $M=X\sq Y$. Then $M$ is decomposable if and only if there exist $k\ge1$, $i\in\Ib$, $j\in\Jb$ and $u\in\Hom_B(V,M_k)$ with $f_{ij1}(u(x_i\otimes y_j))\ne0$. Such a tuple $(k,i,j,u)$ is called a \emph{witness}.
\end{lemmaII}

\begin{proof}
\begin{enumerate}[label=\textup{\arabic*.},leftmargin=1.5em,itemsep=2pt]
\item \emph{The map $\rho$.} Every $\varphi\in\End M$ acts on the multiplicity-free top by scalars, which gives an algebra map $\rho:\End M\to\kk^{\Ib\times\Jb\times\{0,1\}}$. Its kernel consists of maps into $\rad M=\bigcap\ker f_{ije}$, which are nilpotent (as in Proposition~\ref{prop:crit}). So $\ker\rho\subseteq J(\End M)$. The image $\operatorname{im}\rho$ is a subalgebra of some $\kk^s$, hence reduced, hence $\cong\kk^t$ and semisimple. Therefore $J(\End M)=\ker\rho$ and $\End M/J\cong\kk^t$, where $t$ is the number of indecomposable summands of $M$. In particular, $M$ is indecomposable iff $\operatorname{im}\rho=\kk\cdot1$. The coordinate of $\rho(\varphi)$ at $(i,j,e)$ is $f_{ije}(\varphi(1\otimes x_i\otimes y_j))$, because $f_{ije}(1\otimes x_i\otimes y_j)=1$.
\item \emph{Frobenius.} $\varphi$ corresponds to $u=\varphi|_{1\otimes V}\in\Hom_B(V,\Res M)=\bigoplus_k\Hom_B(V,M_k)$ (Lemma~\ref{lem:mackey}). Write $u=\sum_ku_k$. The $u_k$ can be chosen independently.
\item \emph{$k=0$.} Here $u_0\in\End_B(V)=\End X\otimes\End Y$; write $u_0=\sum_s\alpha_s\otimes\beta_s$. Both rings are local. By step~1 applied to $X$ (respectively $Y$), every element of $J(\End X)$ is nilpotent, so it lies in $\ker\rho^X$. Hence $\rho^X(\alpha_s)=a_s\cdot1$ and $\rho^Y(\beta_s)=b_s\cdot1$ are scalars. Since $f_{ije}|_{M_0}=\varepsilon_i\otimes\delta_j$, each $f_{ije}$ contributes the same scalar $c_0=\sum_sa_sb_s$ from $u_0$.
\item \emph{$k\ge1$.} $f_{ij0}$ vanishes on $M_k$: since $a\notin Z_{ij0}$, $W_{Z_{ij0}}\cap W^{(a,b)}=\{\mathrm{id}\}$, while $k\ge1$ forces $w\ne\mathrm{id}$.
\item \emph{Assembling.} Steps 3 and 4 give $\rho(\varphi)_{ij0}=c_0$ and $\rho(\varphi)_{ij1}=c_0+\sum_{k\ge1}f_{ij1}(u_k(x_i\otimes y_j))$. If no witness exists, then $\rho(\varphi)=c_0\cdot1$ for every $\varphi$. If $(k,i,j,u)$ is a witness, take $u_k=u$ and all other $u_{k'}=0$. Then $\rho(\varphi)_{ij0}=0\ne\rho(\varphi)_{ij1}$, so $\rho(\varphi)$ is not scalar.\qedhere
\end{enumerate}
\end{proof}

\noindent\textbf{Region formula.} Let $[\alpha,\beta]$ be the component of $Z_{ij1}$ containing $a$, as a generator interval. Then $\alpha=a-r'$ and $\beta=a+s'$, where $r'$ is the length of the run $a-1,a-2,\dots$ in $D_i$ and $s'$ is the length of the run $1,2,\dots$ in $E_j$. Exactly as in Lemma~\ref{lem:f1f2}(b), for a shuffle $w$:
\[
w\in W_{Z_{ij1}}\iff w\in W_{[\alpha,\beta]}\iff x_1=\dots=x_{\alpha-1}=A\ \text{ and }\ x_{\beta+2}=\dots=x_N=B.
\]

\section{Key Lemma: a witness makes \texorpdfstring{$X$}{X} a summand of a smaller induced module}

\begin{lemmaII}[Key Lemma]\label{II:key}
Assume (H), and let $(k,i,j,u)$ be a witness. Then:
\begin{enumerate}[label=\textup{(\alph*)},itemsep=1pt]
\item $[a-k+1,a-1]\subseteq D_i$ and $[1,k-1]\subseteq E_j$;
\item $X$ is a direct summand of
\[
\widehat L:=\Ind_{H_{a-k}\otimes H_k}^{H_a}\bigl(\Res_{H_{a-k}}X\sq C_{[k-1]}\bigr),
\]
via $H_a$-maps $X\to\widehat L\to X$ whose composite $g$ satisfies $\varepsilon_i\circ g\ne0$. Explicitly $g=\zeta\circ\Psi''$, with $\zeta$ and $\Psi''$ defined in the proof.
\end{enumerate}
Here $H_{a-k}=\langle\pi_1,\dots,\pi_{a-k-1}\rangle$ and $H_k=\langle\pi_{a-k+1},\dots,\pi_{a-1}\rangle$ inside $H_a$.
\end{lemmaII}

\begin{proof}
\begin{enumerate}[label=\textup{Step \arabic*.},leftmargin=3.6em,itemsep=3pt]
\item \emph{$M_k$ as an $H_a$-module.} By Lemma~\ref{lem:cut}(a), the generators $\pi_i$ with $i<a$ change only the left word $\ell$. They act on $V$ by $\pi_p\otimes1$ with $p\le a-k-1$ (case $AA$) or by $1\otimes\pi_q$ with $q\le k-1$ (case $BB$). So for each right word $r$, the space $\widetilde L_r:=\operatorname{span}\{T_{(\ell\mid r)}\otimes v\}$ is $H_a$-stable, and $M_k|_{H_a}=\bigoplus_r\widetilde L_r$. The word model of \S\ref{sec:conv} identifies every $\widetilde L_r$ with
\[
\widetilde L:=\Ind_{H_{a-k}\otimes H_k}^{H_a}\bigl(\Res_{H_{a-k}}X\sq\Res_{H_k}Y\bigr),\qquad T_{(\ell\mid r)}\otimes v\mapsto T_\ell\otimes v,
\]
where $\Res_{H_k}Y$ means restriction to $\langle\pi_1,\dots,\pi_{k-1}\rangle\subseteq H_b$.
\item \emph{The functional on a copy.} By the region formula,
\[
f_{ij1}(T_{(\ell\mid r)}\otimes x\otimes y)=[\ell_1=\dots=\ell_{\alpha-1}=A]\cdot[r\ \ok]\cdot\varepsilon_i(x)\delta_j(y),
\]
where ``$r$ ok'' means $r_{\beta+2-a}=\dots=r_b=B$. Write ``$\ell$ ok'' for $\ell_1=\dots=\ell_{\alpha-1}=A$. On a copy $\widetilde L_r$ with $r$ ok, the functional
\[
\widetilde\lambda(T_\ell\otimes x\otimes y):=[\ell\ \ok]\,\varepsilon_i(x)\delta_j(y)
\]
equals $f_{ij1}|_{\widetilde L_r}$. Since $f_{ij1}$ is $H_N$-linear and $\Res_{H_a}C_{Z_{ij1}}=C_{D_i}$, $\widetilde\lambda$ is $H_a$-linear into $C_{D_i}$. On copies with $r$ not ok, $f_{ij1}$ vanishes.
\item \emph{Choice of copy.} The map $\Psi:X\to M_k$, $x\mapsto u(x\otimes y_j)$, is $H_a$-linear because $u$ is $B$-linear. Write $\Psi=\sum_r\Psi_r$ with $\Psi_r:X\to\widetilde L_r$. Since $f_{ij1}(\Psi(x_i))\ne0$, some copy $r$ (necessarily with $r$ ok) has $\widetilde\lambda(\Psi_r(x_i))\ne0$. Fix it.
\item \emph{Proof of (a).} For $\widetilde\lambda\neq0$ we need a left word whose $k$ letters $B$ all lie in positions $[\alpha,a]$. So $k\le a-\alpha+1=r'+1$, which gives $[a-k+1,a-1]\subseteq[a-r',a-1]\subseteq D_i$.

For the second inclusion, let $E':=E_j\cap[k-1]$. The restriction $\delta_j:\Res_{H_k}Y\to C_{E'}$ is $H_k$-linear, so $\mathrm{id}\otimes\delta_j$ induces a surjective $H_a$-map
\[
\Ind(\mathrm{id}\sq\delta_j):\widetilde L\to\widehat L':=\Ind_{H_{a-k}\otimes H_k}^{H_a}\bigl(\Res_{H_{a-k}}X\sq C_{E'}\bigr),
\]
and $\widetilde\lambda=\widehat\lambda\circ\Ind(\mathrm{id}\sq\delta_j)$, where $\widehat\lambda(T_\ell\otimes x\otimes1):=[\ell\ \ok]\,\varepsilon_i(x)$. Because $\Ind(\mathrm{id}\sq\delta_j)$ is onto and $\widetilde\lambda$ is $H_a$-linear, $\widehat\lambda$ is $H_a$-linear. Now take $z=T_{\mathrm{id}}\otimes x_i\otimes1$, i.e.\ $\ell=A^{a-k}B^k$. Since $\alpha-1\le a-k$ by the first part, $\ell$ is ok and $\widehat\lambda(z)=1$. Apply $\pi_{a-k+q}$ for $q\in[1,k-1]$:
\begin{itemize}[itemsep=1pt]
\item by rule $BB$, $\pi_{a-k+q}z=[q\in E']\,z$;
\item by $H_a$-linearity into $C_{D_i}$, $\widehat\lambda(\pi_{a-k+q}z)=[a-k+q\in D_i]=1$, using the first inclusion.
\end{itemize}
Hence $[q\in E']=1$ for all $q\in[1,k-1]$, so $E'=[k-1]$, i.e.\ $[1,k-1]\subseteq E_j$. In particular $\widehat L'=\widehat L$.
\item \emph{The return map $\zeta$.} Let $w_0^{(k)}$ be the longest element of $W_{[a-k+1,a-1]}$ (the symmetric group on the values $a-k+1,\dots,a$). Define $\theta:\Res_{H_{a-k}}X\otimes C_{[k-1]}\to\Res_{H_{a-k}\otimes H_k}X$ by $\theta(x\otimes1):=\pi_{w_0^{(k)}}x$.
\begin{itemize}[itemsep=1pt]
\item $\theta$ is $H_{a-k}$-linear, because $\pi_{w_0^{(k)}}$ is a word in $\pi_{a-k+1},\dots,\pi_{a-1}$, which commute with $\pi_1,\dots,\pi_{a-k-1}$.
\item $\theta$ is $H_k$-linear for the trivial action $\pi_{a-k+q}\mapsto1$ on $C_{[k-1]}$: $\pi_{a-k+q}\pi_{w_0^{(k)}}=\pi_{w_0^{(k)}}$ by \eqref{eq:mult}, since every $s_{a-k+q}$ is a left descent of $w_0^{(k)}$.
\end{itemize}
Frobenius reciprocity gives the $H_a$-map
\[
\zeta:\widehat L\to X,\qquad \zeta(T_\ell\otimes x\otimes1)=T_\ell\,\pi_{w_0^{(k)}}x.
\]
\item \emph{$\varepsilon_i\circ\zeta=\widehat\lambda$.} Since $\varepsilon_i$ is $H_a$-linear into $C_{D_i}$,
\[
\varepsilon_i(T_\ell\pi_{w_0^{(k)}}x)=\chi_{D_i}(T_\ell)\cdot\chi_{D_i}(\pi_{w_0^{(k)}})\cdot\varepsilon_i(x).
\]
The middle factor is $[w_0^{(k)}\in W_{D_i}]=1$ by (a). For the first factor, $\ell$ is an $(a-k,k)$-shuffle in $S_a$, so its right descents lie in $\{a-k\}$. Write $W_{D_i}$ as the commuting product of $W_K$ over the components $K$ of $D_i$. The component containing $a-1$ is the generator interval $[\alpha,a-1]$ (empty if $r'=0$). All other components lie in $[1,\alpha-2]$, and $\alpha-2\le a-k-1$ because $k\le r'+1$. As in Lemma~\ref{lem:f1f2}(b), a factor in a component not containing $a-k$ would contribute a right descent $\ne a-k$, so
\[
\ell\in W_{D_i}\iff\ell\in W_{[\alpha,a-1]}\qquad(\text{generator interval } \{s_\alpha,\dots,s_{a-1}\}).
\]
In the edge case $k=r'+1$ we have $a-k=\alpha-1\notin D_i$, both sides force $\ell=\mathrm{id}$, and $\ell=A^{a-k}B^k$ is exactly the unique ok word, so the equivalence still holds. $W_{[\alpha,a-1]}$ is the symmetric group on the values $\{\alpha,\dots,a\}$. So, as in the proof of Lemma~\ref{lem:f1f2}(b), $\ell\in W_{[\alpha,a-1]}$ iff $\ell$ fixes the values $1,\dots,\alpha-1$, iff $\ell_1=\dots=\ell_{\alpha-1}=A$, iff ``$\ell$ ok''. Hence $\varepsilon_i\circ\zeta=\widehat\lambda$.
\item \emph{Conclusion.} Put $\Psi'':=\Ind(\mathrm{id}\sq\delta_j)\circ\Psi_r:X\to\widehat L$ and $g:=\zeta\circ\Psi''\in\End_{H_a}X$. Then
\[
\varepsilon_i(g(x_i))=\widehat\lambda(\Psi''(x_i))=\widetilde\lambda(\Psi_r(x_i))\ne0.
\]
So $g\notin\ker\rho^X=J(\End X)$, and $g$ is a unit of the local ring $\End X$. Then $X\xrightarrow{\Psi''}\widehat L\xrightarrow{g^{-1}\zeta}X$ is the identity, and $X$ is a summand of $\widehat L$.\qedhere
\end{enumerate}
\end{proof}

\section{Shape input}

\begin{lemmaII}[connected shapes; KLO Lemma 5.3(1), 5.4(1), Thm.~3.9]\label{II:conn}
Let $C$ be a connected skew shape, and let $T_0$ be its row-superstandard tableau: rows top to bottom, each filled left to right with consecutive values. Then $X_C$ has simple top $C_{D_C}$, where $D_C=\{i: i,i+1\text{ lie in the same row of }T_0\}$, and its top functional is $\varepsilon_C(T)=[T=T_0]$.
\end{lemmaII}

\begin{proof}
\emph{Simple top.} By KLO Lemma 5.4(1), the projective cover of $X_C$ is $P_{\alpha_{\mathrm{proj}}(C)}$, which is indecomposable since $\alpha_{\mathrm{proj}}(C)$ is an ordinary composition. Its label is $D_C$, the set of $i$ with $\pi_iT_0=T_0$. Indeed, if $i$ ends row $r$ and $i+1$ starts row $r+1$, connectedness puts $i+1$ weakly left of $i$, so $\pi_iT_0\in\{0,s_iT_0\}$. This was cross-checked against the Wake 289 sweep data.

\emph{$\sigma$ explicitly.} In the interval model (KLO Thm.~3.9, Lemma~3.2), $X_C$ is a left weak Bruhat interval module $B(\sigma,\rho)$ with $T_0\mapsto\sigma:=\mathrm{read}_{\tau_0}(T_0)$. This reading takes rows right to left, top row first. So $\sigma$ concatenates \emph{decreasing} runs on consecutive blocks of values, one block per row. Hence $\Des_L(\sigma)=D_C$, and $\sigma$ is the longest element of $W_{D_C}$.

\emph{Reachability.} Every $\gamma\ne\sigma$ in $[\sigma,\rho]_L$ is $u\sigma$ for some $u\ne e$ with $\ell(u\sigma)=\ell(u)+\ell(\sigma)$. A reduced word of $u$ gives a saturated ascending chain inside $[\sigma,\gamma]_L\subseteq[\sigma,\rho]_L$. So $\gamma=\pi_u\sigma$, and $\varepsilon(\gamma)=[u\in W_{D_C}]\,\varepsilon(\sigma)$.

\emph{Vanishing.} If $u\in W_{D_C}$ and $u\neq e$, then $u\sigma\in W_{D_C}$ would have length $>\ell(\sigma)$, but $\ell(\sigma)$ is the maximum length in $W_{D_C}$. So $\varepsilon(\gamma)=0$ for all $\gamma\ne\sigma$.
\end{proof}

\begin{corII}\label{II:mf}
Every $X_{\lambda/\mu}$ with $m$ components has a multiplicity-free top of size $2^{m-1}$. This follows from Lemmas~\ref{II:W}, \ref{II:top} and~\ref{II:conn} by induction on the number of components.
\end{corII}

\begin{lemmaII}[support]\label{II:support}
Let $X=X_{C_1\star\cdots\star C_h}$ with $|X|=a$ and $C_h$ not a ribbon, and let $\ell_{\mathrm{last}}$ be the length of the last row of $C_h$ (the bottom row of the shape). Let $\varepsilon$ be a top functional of $X$ with label $D\supseteq[a-k+1,a-1]$, where $k\ge1$. Then:
\begin{enumerate}[label=\textup{(\roman*)},itemsep=1pt]
\item $k\le\ell_{\mathrm{last}}$;
\item if $\varepsilon(T)\ne0$, then the values $a-k+1,\dots,a$ of $T$ occupy the last $k$ cells of the bottom row.
\end{enumerate}
\end{lemmaII}

\begin{proof}
Put $a':=a-|C_h|$.

\emph{Bound on $k$.} By Lemma~\ref{II:top} (iterated through Lemma~\ref{II:W}), $D\cap[a'+1,a-1]=a'+D_{C_h}$. The final run of $D_{C_h}$ consists of the pairs inside the last row of $T_0^{C_h}$. It has length $\ell_{\mathrm{last}}-1$ and ends at $|C_h|-1$. It is preceded by a non-element, since $C_h$ has at least two rows. So the run of $D$ ending at $a-1$ has length exactly $\ell_{\mathrm{last}}-1$, and $[a-k+1,a-1]\subseteq D$ forces $k-1\le\ell_{\mathrm{last}}-1$.

\emph{$h=1$.} By Lemma~\ref{II:conn}, only $T=T_0$ qualifies, and $T_0$ has $a-k+1,\dots,a$ at the end of the last row.

\emph{$h\ge2$.} Write $X=X'\sq X_{C_h}$ by Lemma~\ref{II:W}, with $X'=X_{C_1\star\cdots\star C_{h-1}}$, $A$ marking $X'$-values and $B$ marking $C_h$-values. By Lemma~\ref{II:top}, $\varepsilon=f_{i'e}$ with $\varepsilon(T_w\otimes x'\otimes y)=[w\in W_Z]\,\varepsilon'(x')\,\varepsilon_{C_h}(y)$.
\begin{itemize}[itemsep=1pt]
\item $\varepsilon_{C_h}(y)\ne0$ forces $y=T_0^{C_h}$ (Lemma~\ref{II:conn}).
\item If $e=0$, then $w=\mathrm{id}$. If $e=1$, then $w\in W_{[\alpha,\beta]}$ with $\beta=a'+\ell_{\mathrm{first}}(C_h)-1$ (the run $1,2,\dots$ in $D_{C_h}$ has length $\ell_{\mathrm{first}}-1$). So $w$ fixes every value $>a'+\ell_{\mathrm{first}}(C_h)$.
\end{itemize}
Since $C_h$ has at least two rows, $|C_h|\ge\ell_{\mathrm{first}}+\ell_{\mathrm{last}}\ge\ell_{\mathrm{first}}+k$. So $a-k+1>a'+\ell_{\mathrm{first}}$, and the values $a-k+1,\dots,a$ are fixed by $w$. They therefore sit in their $T_0^{C_h}$ positions, namely the last $k$ cells of the last row.
\end{proof}

\begin{remarkII}[non-projectivity; not used below]\label{II:nonproj}
Let $C$ be connected and not a ribbon. Then $X_C$ is not projective.

\emph{Proof.} $X_C$ has simple top, with projective cover $P_\alpha$, $\alpha=\alpha_{\mathrm{proj}}(C)$. Let $R_\alpha$ be the ribbon with the same row lengths; then $P_\alpha\cong X_{R_\alpha}$ (KLO Rem.~A.1(2)). Identify cells of $C$ and of $R_\alpha$ by (row, index in row). The only $R_\alpha$-relation between rows $r$ and $r+1$ is $(r,1)<(r+1,\text{last})$. In $C$, $(r,1)$ lies above a cell of row $r+1$, and that cell is $\le(r+1,\text{last})$. So every relation of $R_\alpha$ holds in $C$, and $\SYT(C)\subseteq\SYT(R_\alpha)$. If rows $r$ and $r+1$ of $C$ overlap in $o\ge2$ columns, $C$ forces $(r,o)<(r+1,\text{last})$. In $R_\alpha$ these two cells are incomparable, since the up-set of $(r,o)$ is the rest of row $r$. So some SYT of $R_\alpha$ violates $C$, $\dim X_C<\dim P_\alpha$, and $X_C\not\cong P_\alpha$. \qed
\end{remarkII}

\section{Proof of Theorem~\ref{thm:B} and Corollary~\ref{cor:C}}

We induct on $n=|\lambda/\mu|$. For $m=1$, KLO Prop.~7.3(1) shows that connected shapes give indecomposable modules (simple top). Let $m\ge2$ and assume Theorem~\ref{thm:B} for all sizes $<n$. Corollary~\ref{cor:C} at sizes $<n$ then follows from Theorem~\ref{thm:B} there, by the argument at the end of this section.

\begin{enumerate}[label=\textup{Step \arabic*.},leftmargin=3.6em,itemsep=3pt]
\item \emph{Choose a split.} Pick $h<m$ with $C_h$ not a ribbon. If there is none, then $C_1,\dots,C_{m-1}$ are all ribbons, and the hypothesis forces $m=2$ with $C_1$ a ribbon and $C_2$ not. In that case Lemma~\ref{II:rot} replaces $\lambda/\mu$ by its rotation $C_2^\circ\star C_1^\circ$, and we take $h=1$.
\item \emph{Set up (H).} Put $X=X_{C_1\star\cdots\star C_h}$ and $Y=X_{C_{h+1}\star\cdots\star C_m}$, with $|X|=a$ and $|Y|=b$, both $<n$. Each factor either inherits the hypothesis of Theorem~\ref{thm:B} or is a single component, so both are indecomposable by induction or by KLO Prop.~7.3(1). Their tops are multiplicity-free by Corollary~\ref{II:mf}. So (H) holds, and $X_{\lambda/\mu}\cong X\sq Y$ by Lemma~\ref{II:W}.
\item \emph{Suppose $X\sq Y$ is decomposable.} Lemma~\ref{II:crit} gives a witness $(k,i,j,u)$. Lemma~\ref{II:key}(a) gives $D_i\supseteq[a-k+1,a-1]$, and then Lemma~\ref{II:support}(i) gives $k\le\ell_{\mathrm{last}}(C_h)$. By Lemma~\ref{II:key}(b), $X$ is a summand of $\widehat L$, via $g=\zeta\circ\Psi''$ with $\varepsilon_i\circ\zeta=\widehat\lambda$.
\item \emph{Restrict to one piece.} Write $\Res_{H_{a-k}}X=\bigoplus_\nu Q_\nu$. Here $\nu$ runs over the order ideals of size $a-k$ of the shape of $X$, and $Q_\nu:=\operatorname{span}\{T:T^{-1}([1,a-k])=\nu\}$. This is a decomposition because $\pi_i$ with $i<a-k$ only moves values $\le a-k$. Moreover $Q_\nu\cong X_\nu^{\oplus f^{\mathrm{sh}(X)/\nu}}$, where $f^{\mathrm{sh}(X)/\nu}$ is the number of SYT of the complementary skew shape.
\begin{itemize}[itemsep=1pt]
\item Accordingly $\widehat L=\bigoplus_\nu\widehat L_\nu$ with $\widehat L_\nu:=\Ind(Q_\nu\sq C_{[k-1]})$, and $g=\sum_\nu g_\nu$ with $g_\nu:=\zeta\,\iota_\nu\,p_\nu\,\Psi''$, where $p_\nu$ and $\iota_\nu$ are the projection and inclusion.
\item $\End X$ is local, so the non-units form an ideal, and some $g_\nu$ is a unit. Then $\varepsilon_i\circ g_\nu\ne0$, so $\widehat\lambda=\varepsilon_i\circ\zeta$ is nonzero on $\widehat L_\nu$. Since $\widehat\lambda(T_\ell\otimes x\otimes1)=[\ell\ \ok]\varepsilon_i(x)$, this gives an SYT $T\in Q_\nu$ with $\varepsilon_i(T)\ne0$.
\item By Lemma~\ref{II:support}(ii), this forces $\nu=\nu_0$, the shape of $X$ minus the last $k$ cells of the bottom row. This is a valid order ideal because $k\le\ell_{\mathrm{last}}$, and $f^{\mathrm{sh}(X)/\nu_0}=1$.
\item Hence $X$ is a direct summand of $\widehat L_{\nu_0}\cong X_{\nu_0}\sq X_{(k)}\cong X_{S_0}$, where
\[
S_0:=C_1\star\cdots\star C_{h-1}\star C_h^-\star(k).
\]
This uses Lemma~\ref{II:W} and $C_{[k-1]}=X_{(k)}$, the one-row module. (Empty columns created by deleting cells may be compressed away; the module depends only on the relative row and column order of the cells.) $C_h^-$ may itself split as $C_h'\star(\text{row segment})$.
\end{itemize}
\item \emph{Contradiction.} $|S_0|=a<n$, so Corollary~\ref{cor:C} at size $a$ gives $X_{S_0}\cong\bigoplus_SX_S$ with every $X_S$ indecomposable. In every such $S$ the last component is a ribbon, because the final maximal run of ribbons of $S_0$ contains $(k)$ and is merged into a single ribbon. By Krull--Schmidt, $X\cong X_S$ for one of these $S$. By KLO Thm.~5.5 ($X_{\lambda/\mu}\cong X_{\nu/\kappa}$ iff $\lambda/\mu=\nu/\kappa$), $S=C_1\star\cdots\star C_h$. But the last component of $C_1\star\cdots\star C_h$ is $C_h$, which is not a ribbon. This contradiction shows that $X\sq Y$ is indecomposable.\qed
\end{enumerate}

\noindent\textbf{Proof of Corollary~\ref{cor:C}.} By Lemma~\ref{II:W} and associativity of $\sq$, $X_{\lambda/\mu}\cong X_{C_1}\sq\cdots\sq X_{C_m}$. Ribbon modules are projective, $X_{R_\gamma}\cong P_\gamma$ (KLO Rem.~A.1(2)), and
\[
X_R\sq X_{R'}\cong P_\gamma\sq P_{\gamma'}\cong P_{\gamma\cdot\gamma'}\oplus P_{\gamma\odot\gamma'}\cong X_{R''}\oplus X_{R'''}
\]
(KLO Lemma 5.1; see Remark~\ref{II:huang}). Merge each maximal run of ribbons in this way, using distributivity of $\sq$ over $\oplus$. Every resulting shape $S$ has no two consecutive ribbon components, so each $X_S$ is indecomposable by Theorem~\ref{thm:B}, or by KLO Prop.~7.3(1) if $S$ is connected.\qed

\begin{remarkII}[Huang's lemma is used second-hand]\label{II:huang}
The identity $P_\gamma\sq P_{\gamma'}\cong P_{\gamma\cdot\gamma'}\oplus P_{\gamma\odot\gamma'}$ is KLO Lemma 5.1, which KLO attribute to Huang (their reference [18], Thm.~3.3). \textbf{I have used it second-hand through KLO and have not consulted Huang's paper.} The identity can, however, be re-derived here, as in the converse of Theorem~\ref{thm:A}. $P_\gamma\otimes P_{\gamma'}$ is a projective $B$-module, so $P_\gamma\sq P_{\gamma'}$ is projective. By Lemma~\ref{II:top} its top is $C_{D_{\gamma\odot\gamma'}}\oplus C_{D_{\gamma\cdot\gamma'}}$, multiplicity-free. Hence it is $P_{\gamma\cdot\gamma'}\oplus P_{\gamma\odot\gamma'}$. The identification of these projectives with the two glued ribbon modules is KLO Rem.~A.1(2), which I read first-hand.
\end{remarkII}

\section{Remarks on Part II}

\begin{remarkII}[ribbon at the end of $X$]
If $C_h$ is a ribbon, Lemma~\ref{II:key} gives no contradiction. In that case $X=X_{\cdots\star R}$ genuinely \emph{is} a summand of $X_{\cdots\star R^-\star(k)}$, because one merge pattern reproduces $R$. This is exactly the disconnected-ribbon splitting of KLO Prop.~7.3(2).
\end{remarkII}

\begin{remarkII}[(G*) is false]
The tempting general statement (G*) is: ``Ind of two indecomposables is indecomposable unless both are projective.'' It would prove Theorem~\ref{thm:B} by associativity, but it is \textbf{false}: \texttt{gstar.py} finds 56 counterexamples among weak-Bruhat interval modules with $a+b\le7$. For example, $C_{(1)}\sq Y$ decomposes for some non-projective indecomposable interval modules $Y$ of $H_4(0)$. Shape modules avoid this failure through two ingredients. The support lemma (Lemma~\ref{II:support}) pins the summand to one restriction piece. The classification theorem (KLO Thm.~5.5) then makes ``summand of $X_{S_0}$'' decidable by comparing last components.
\end{remarkII}

\begin{remarkII}[relation to Part I]
Q72(i) has the same Lemma~\ref{II:crit} $+$ Lemma~\ref{II:key} skeleton with one-dimensional factors. There, $\widehat L$ becomes $L_k$, and the contradiction comes from Norton (Fact~\ref{fact:N}) instead of KLO Thm.~5.5.
\end{remarkII}

\section{Cited results}\label{sec:cited}
\begin{itemize}[itemsep=2pt]
\item \textbf{Part I:} Norton's classification of simples and $\dim P_D$ (Fact~\ref{fact:N}), and freeness of $H_N$ over a parabolic subalgebra. Deodhar's lemma is re-proved inline in the only case needed (Lemma~\ref{lem:rules}(iii)).
\item \textbf{Part II.} All of the following are from KLO arXiv:2310.20571 (\S5, \S7.2 and App.~A were read first-hand in PROVE 290):
\begin{itemize}[itemsep=1pt]
\item Thm.~3.9 / Lemma 5.3(1): the interval basis with $\sigma=\mathrm{read}(T_0)$, used only in Lemma~\ref{II:conn};
\item Lemma 5.4(1): the projective cover $P_{\alpha_{\mathrm{proj}}}$;
\item Rem.~A.1(2): ribbon modules are projective;
\item Lemma 5.1 (Huang): second-hand, see Remark~\ref{II:huang};
\item Thm.~5.5: $X_{\lambda/\mu}\cong X_{\nu/\kappa}$ iff $\lambda/\mu=\nu/\kappa$;
\item Prop.~7.3(1): connected shapes give indecomposable modules.
\end{itemize}
Lemma~\ref{II:W} is proved here directly; it is KLO Prop.~A.2(1) / JKLO arXiv:2105.14169 Lemma~4.
\end{itemize}

\section{Computational evidence}
\begin{itemize}[itemsep=2pt]
\item \textbf{Q72(i).} \texttt{proofs/prove289/check.py 8} verifies every intermediate claim of Part~I for $N\le8$: (a) $f_1$ and $f_2$ from \eqref{eq:f2} are module maps; (b) (Crit) separation $\iff$ both projective; (c) the per-$k$ criterion, computed inside the small modules $L_k,R_k$, $\iff$ both projective. All 769 cases $(a,D_I,D_J)$ with $N\le8$ pass, with 0 mismatches.

Independent prior data (Wake 288, computed mod $p=1000003$, with exact relation checks for $N\le4$): every non-both-projective $M_{I,J}$ is indecomposable for $N=4,\dots,8$ (4/20/64/172/424 cases, 0 exceptions), and the both-projective cases split into exactly two summands ($N\le6$, 31 cases).
\item \textbf{Q73.} \texttt{chain2.py} runs every split $X\sq Y$ with indecomposable factors in the explicit word model:
\begin{itemize}[itemsep=1pt]
\item coverage: all 2-component shapes with $n\le7$ (336 splits) and all 3-component shapes with $6\le n\le7$ under both splits (36 splits, including factors with a 2-dimensional top);
\item ``witness exists'' $\iff$ ``decomposable'' $\iff$ ``contains a disconnected ribbon'', in agreement with the independent Wake 289 End-algebra sweep;
\item $f_{ij1}$ is a module map, Lemma~\ref{II:key}(a) holds for every witness, and no witness exists when the last component of $X$ is not a ribbon;
\item 0 failures.
\end{itemize}
\texttt{chain.py kl 5}: in all 54 witness cases with $n\le5$, $X$ is a direct summand of $\widehat L$ (a random composite $X\to\widehat L\to X$ is invertible). Wake 289: Conjecture 7.5 holds for every shape with $n\le8$ (304 hypothesis shapes). KLO themselves checked $|P|\le6$.
\end{itemize}

\section{Status and owner risk}
\noindent\textbf{Q72(i).} Owner gate (Wake 288): NOT FOUND. Jung--Kim--Lee--Oh (arXiv:2105.14169, \S5) pose the open problem of decomposing weak Bruhat interval modules $B(\sigma,\rho)$ into indecomposables. By their product rule, $M_{I,J}$ is a single interval module, so Q72(i) is a special case of that problem. Tewari--van Willigenburg (Kronecker $0$-Hecke) is not yet checked.

\noindent\textbf{Q73.} Owner risk is \textbf{open}. The Wake 289 owner gate and the Browse 242 citation sweep found no proof among papers citing KLO, and I estimate about 85\% that the conjecture is still open in print. The main risk is the Kim--Lee--Oh school, which may have a proof in preparation. Defant--Searles, arXiv:2404.04961, has \textbf{not yet been checked} and must be read before any public claim.

\noindent\textbf{What I want from you.} A referee-style check of both proofs, especially Lemma~\ref{II:key} (the Key Lemma, = Lemma 3 of the PROVE 290 record: copy decomposition, factorisation $\widetilde\lambda=\widehat\lambda\circ\Ind(\mathrm{id}\sq\delta_j)$, return map $\zeta$ and the identity $\varepsilon_i\circ\zeta=\widehat\lambda$) and Lemma~\ref{II:support} (support, = Lemma 5 of the record, case $h\ge2$). If you know of a source for either theorem, that is equally valuable.

\end{document}
